\begin{frame}{Conjugação Topológica}
  \begin{block}{Definição}
    Dizemos que uma função $h:X\to Y$ é um \textbf{homeomorfismo} se
    \begin{enumerate}
      \item $h$ é bijetiva.
      \item $h$ é contínua.
      \item $h^{-1}$ é contínua. 
    \end{enumerate}
  \end{block}
  Exemplo: Seja $h:\mathbb{R}\to\mathbb{R}$ dada por $h(x)=\sqrt[3]{x}$.\\
  Imagem
\end{frame}

\begin{frame}{}
  \begin{block}{Definição}
    Considere os dois seguintes sistemas autónomos
    \begin{equation}
      \begin{cases}
        x^{\prime}=f(x) \text{,} \\
        y^{\prime}=g(y).
      \end{cases}
    \end{equation}
    com seus respectivos fluxos $\phi_t(x)$ e $\varphi_t(y)$.\\
    Dizemos que $h$ é uma \textbf{conjugação topológica} dos sistemas autónomos se $h$ é um homeomorfismo tal que
    \begin{align*}
      h(\phi_t(x))=\varphi_t(h(x)),\quad\text{para todo }t\in\mathbb{R}.
    \end{align*}
    Grafo Algebristas :p.
  \end{block}
\end{frame}

\begin{frame}{}
  \begin{block}{Observação}
    Seja $h$ conjugação topológica, então
    \begin{align*}
      h\left( \mathcal{O}_{\phi}(x) \right)=&h\left( \left\{ \phi_t(x):t\in\mathbb{R} \right\} \right)\\
      =&\left\{ h(\phi_t(x)):t\in\mathbb{R} \right\}\\
      =&\left\{ \varphi_t(h(x)):t\in\mathbb{R} \right\}\\
      =&\mathcal{O}_{\varphi}(h(x)). 
    \end{align*}
    Pelo que podemos dizer que $h$ leva órbitas de $\phi$ em órbitas de $\varphi$, i.e. $h$ leva a dinâmica de um sistema ao outro.\\
    Em particular, $h$ leva singularidades em singularidades. 
  \end{block}
\end{frame}
